% ===========================================================
%  第 4 章 矩阵的秩 —— 北大《高等代数》第五版【深化】提取
%  来源：PDF 90–103 = 印刷 78–91
%        印刷 78–83 = 第三章 §3 线性相关性
%        印刷 84–91 = 第三章 §4 矩阵的秩
%  对应 OUTLINE.md §5 深化清单 4.1 / 4.2（4.3、4.4 不在本页范围，见报告）
%  注意：本文件头部注释不算块，提取脚本只认行首的 \begin{fdeep}
% ===========================================================

% ---------- 印刷 81（§3）----------
\begin{fdeep}{线性相关性的方程组判定与“延长”性质（北大 三.3，印刷 81）}
向量组 $\alpha_1,\dots,\alpha_s$（其中 $\alpha_i=(a_{i1},a_{i2},\dots,a_{in})$）是否线性相关，
等价于齐次方程组
\fml{x_1\alpha_1+x_2\alpha_2+\cdots+x_s\alpha_s=0}
有无非零解；按分量写出来就是
\[
\begin{cases}
a_{11}x_1+a_{21}x_2+\cdots+a_{s1}x_s=0,\\
a_{12}x_1+a_{22}x_2+\cdots+a_{s2}x_s=0,\\
\cdots\cdots\cdots\cdots\\
a_{1n}x_1+a_{2n}x_2+\cdots+a_{sn}x_s=0.
\end{cases}
\]
于是有\textbf{向量组线性相关的充要条件是这个齐次方程组有非零解}，
\textbf{线性无关的充要条件是它只有零解}。
由此还得出一条常用性质（“延长”）：若向量组 $\alpha_1,\dots,\alpha_s$ 线性无关，
则在每一个向量上添一个分量所得的向量组也线性无关。理由是添分量相当于给上面的
方程组多添一个方程，方程变多解只会变少，原方程组只有零解，新方程组自然也只有零解。
【反哺点】服务考纲〔向量〕“理解向量组线性相关、线性无关的概念”与〔线性方程组〕“理解齐次线性方程组有非零解的充分必要条件”。它把“相关／无关”从文字定义变成“解方程组”，是秩与方程组之间的第一座桥：矩阵 $A$ 的列向量组线性相关 $\iff Ax=0$ 有非零解 $\iff r(A)<$ 列数。“延长性质”给出免算技巧：短向量组无关 $\Rightarrow$ 添分量后仍无关（逆否：长向量组相关 $\Rightarrow$ 去掉分量后仍相关），凡题目里“给向量补一个分量证明仍无关”的题，一句“方程多了解不会变多”即可说清。
\end{fdeep}

% ---------- 印刷 82（§3 定理 2）----------
\begin{fdeep}{替换定理：向量组个数的比较定理（北大 三.3 定理 2 及推论，印刷 82）}
设 $\alpha_1,\dots,\alpha_r$ 可由 $\beta_1,\dots,\beta_s$ 线性表出。若 $r>s$，则 $\alpha_1,\dots,\alpha_r$ 必线性相关。
\fml{x_1\alpha_1+\cdots+x_r\alpha_r=\sum_{i=1}^{r}x_i\sum_{j=1}^{s}t_{ji}\beta_j=\sum_{j=1}^{s}\Bigl(\sum_{i=1}^{r}t_{ji}x_i\Bigr)\beta_j}
证明思路：设 $\alpha_i=\sum_{j=1}^{s}t_{ji}\beta_j$，作上面的线性组合。要得到零组合，
只需让 $\beta_1,\dots,\beta_s$ 的系数全为 $0$，那是 $s$ 个方程、$r$ 个未知量的齐次方程组；
由 $r>s$（未知量个数大于方程个数）它必有非零解，于是 $\alpha$ 组线性相关。
\textbf{全部力量只来自一句话：未知量多于方程的齐次方程组必有非零解。}
由此立得三条推论：①若 $\alpha$ 组可由 $\beta$ 组线性表出，且 $\alpha$ 组线性无关，则 $r\le s$；
②任意 $n+1$ 个 $n$ 维向量必线性相关；③两个线性无关且等价的向量组必含有相同个数的向量。
【反哺点】服务考纲〔向量〕“理解向量组的秩……的概念”。推论 ① 是“比较向量个数”的唯一常用武器：要证 $m\le k$，就把 $m$ 个线性无关的向量用 $k$ 个向量表出，再用反证法；推论 ② 直接给出“$n$ 维空间里向量一多必相关”，是判断含参向量组相关性的快捷下界；推论 ③ 是“向量组的秩”能够良定义的依据，也是下一节“行秩 $=$ 列秩”证明中反复比较个数的底层定理。
\end{fdeep}

% ---------- 印刷 83（§3 定义 13、14、定理 3）----------
\begin{fdeep}{极大线性无关组与向量组的秩（北大 三.3 定义 13、14 与定理 3，印刷 83）}
向量组的一个部分组若本身线性无关，且从原向量组中任添一个向量进来所得的部分组都线性相关，
就称为原向量组的一个\textbf{极大线性无关组}。两条要点：
\begin{itemize}[leftmargin=2.2em,itemsep=0.22em]
\item 任一极大线性无关组都与向量组本身\textbf{等价}。部分组表出整体是显然的；反过来，
      对组中任一 $\alpha_i$，由极大性 $\alpha_1,\dots,\alpha_s,\alpha_i$ 线性相关，即有不全为零的
      $k_1,\dots,k_s,l$ 使 $k_1\alpha_1+\cdots+k_s\alpha_s+l\alpha_i=0$；若 $l=0$，
      则 $\alpha_1,\dots,\alpha_s$ 线性相关，与假设矛盾，故 $l\ne0$，于是
      $\alpha_i=-\frac{k_1}{l}\alpha_1-\cdots-\frac{k_s}{l}\alpha_s$ 可由部分组线性表出。
\item 既然各个极大无关组都彼此等价，而它们又都线性无关，由替换定理的推论 ③，
      \textbf{一个向量组的任意两个极大线性无关组都含有相同个数的向量}。
\end{itemize}
于是可以放心地\textbf{定义}：向量组的极大线性无关组所含向量的个数，称为这个向量组的\textbf{秩}。
由定义立得：向量组线性无关的充要条件是它的秩等于它所含向量的个数。
【反哺点】服务考纲〔向量〕“理解向量组的秩、极大线性无关组的概念”。秩之所以是一个“确定”的数（可以放心去算），靠的正是“极大无关组彼此等价 $+$ 等价无关组个数相同”这两步，而不是规定。考场上“求秩”“求极大无关组”“判断相关性”三件事本质是同一个操作（化阶梯形、数非零行、取主元列）——先记住这条，再去看第 4 章的矩阵秩。
\end{fdeep}

% ---------- 印刷 84（§3 末 + §4 定义 15 前）----------
\begin{fdeep}{线性表出与方程组的同解（北大 三.3 末，印刷 84）}
把每个方程看成一个“行增广向量”：方程组 $A_1,\dots,A_s$ 对应 $\alpha_i=(a_{i1},\dots,a_{in},d_i)$，
另一个方程 $B$ 对应 $\beta=(b_1,\dots,b_n,d)$。则 $\beta$ 是 $\alpha_1,\dots,\alpha_s$ 的线性组合，
当且仅当方程 $B$ 是方程 $A_1,\dots,A_s$ 的线性组合；而这时 $\alpha$ 组的解一定满足 $B$。
进一步：若 $\beta_1,\dots,\beta_r$ 可经 $\alpha_1,\dots,\alpha_s$ 线性表出，则方程组 $A$ 的解都是方程组 $B$ 的解；
若两个行向量组等价，则两个方程组\textbf{同解}。
【反哺点】服务考纲〔线性方程组〕“会用初等行变换求解线性方程组”。这正是“解方程组可以做初等行变换”的合法性根据：初等行变换前后的行向量组等价，所以新方程组与原方程组同解。更重要的推论（后面推论 3 用它）：行变换可逆，故它保持列向量之间的线性关系（$Ax=0$ 与 $PAx=0$ 同解），所以求 $A$ 的\emph{列}向量组的极大无关组时只许做行变换；同理，求\emph{行}向量组的极大无关组时只许做列变换。考生最常见的错误就是在求列向量组的极大无关组时顺手做了列变换。
\end{fdeep}

% ---------- 印刷 84–85（§4 定义 16、17）----------
\begin{fdeep}{秩的子式定义怎么用：$k$ 阶子式与最高阶非零子式（北大 三.4 定义 16、17，印刷 85）}
在 $s\times n$ 矩阵 $A$ 中任取 $k$ 行与 $k$ 列，位于这些行、列交点上的 $k^2$ 个元素按原次序组成的
$k$ 阶行列式，称为 $A$ 的一个 $k$ \textbf{阶子式}（$k\le\min(s,n)$）。
\textbf{定义}：$A$ 中最高阶非零子式的阶数称为 $A$ 的秩，零矩阵的秩规定为 $0$。
用起来更方便的是它的等价说法——$r(A)=r$ 意味着：
\begin{itemize}[leftmargin=2.2em,itemsep=0.22em]
\item 存在一个 $r$ 阶子式不为零；
\item 所有 $r+1$ 阶子式全为零；且由行列式按行（列）展开，$r+2$ 阶、$r+3$ 阶以至更高阶子式也全为零
      （它们都按 $r+1$ 阶子式展开，系数全为 $0$）。
\end{itemize}
\textbf{操作要点}：要证 $r(A)\ge k$，只需\textbf{亮出一个} $k$ 阶非零子式；
要证 $r(A)\le k$，则必须证明\textbf{所有} $k+1$ 阶子式都等于零——只算一个子式是不够的。
【反哺点】服务考纲〔矩阵〕“理解矩阵的秩的概念”。这是四个等价定义中唯一能直接用于\emph{证明}的定义：秩不等式的题一律靠它——“$\ge$”构子式，“$\le$”说全为零。与之配套的还有一条常用技巧：$r(A)\ge k$ 时不必真的算出子式的值，只要说明某 $k$ 行 $k$ 列的行列式非零即可（例如三角块、范德蒙德型、含 $|B|\ne0$ 的分块）。算秩时则一律改用初等变换，子式定义只用来讲理、不用来算。
\end{fdeep}

% ---------- 印刷 86–88（§4 定理 4）----------
\begin{fdeep}{行秩 $=$ 列秩：北大的子式证法（北大 三.4 定理 4，印刷 86--88）}
\textbf{定理}：矩阵 $A$ 的秩 $=$ 列秩 $=$ 行秩（行与列的情形完全对称，只证一半即可）。
设 $A$ 为 $s\times n$ 矩阵，$r(A)=r$，证明按五步走：
\begin{enumerate}[leftmargin=2.2em,itemsep=0.22em]
\item 由子式定义，存在一个 $r$ 阶非零子式 $d$；必要时交换列的次序，可设 $d$ 位于前 $r$ 列、
      第 $i_1,\dots,i_r$ 行上（换列前后两矩阵的列向量组是一样的，有相同的列秩）。
\item $d$ 的 $r$ 个列向量线性无关：因 $d\ne0$，以 $d$ 为系数矩阵的齐次方程组只有零解（克拉默法则）。
\item $A$ 的前 $r$ 个列向量是 $d$ 的 $r$ 个列向量各添上 $s-r$ 个分量得到的，由“延长性质”仍线性无关。
\item 再证 $A$ 的任一列 $\alpha_j$ 都可由前 $r$ 列线性表出：以 $d$ 为系数矩阵、以 $\alpha_j$ 在
      $i_1,\dots,i_r$ 行上的分量为右端的方程组有唯一解 $l_1,\dots,l_r$；作初等列变换把第 $j$ 列换成
      $\alpha_j-l_1\alpha_1-\cdots-l_r\alpha_r$，则它在 $i_1,\dots,i_r$ 行上的分量全为 $0$。
\item 最后用“$r+1$ 阶子式为零”收拾剩下的分量：取 $i_1,\dots,i_r$ 与任一 $i$ 行、前 $r$ 列与第 $j$ 列
      构成的 $r+1$ 阶子式，它必为零；把第 $i$ 行调到第 $1$ 行、第 $j$ 列调到第 $1$ 列（子式仍为零），
      再作列变换消元，即得该子式 $=b_{ij}d=0$，由 $d\ne0$ 得 $b_{ij}=0$。
      于是 $\alpha_j=l_1\alpha_1+\cdots+l_r\alpha_r$。
\end{enumerate}
所以前 $r$ 列是 $A$ 的列向量组的一个极大线性无关组，$A$ 的列秩 $=r=A$ 的秩；行秩同理。
由证明过程还可直接看出：\textbf{秩为 $r$ 的矩阵中，$r$ 阶非零子式所在的那 $r$ 个行（列）向量，
就是行（列）向量组的一个极大线性无关组}。
【反哺点】服务考纲〔向量〕“理解矩阵的秩与其行（列）向量组的秩之间的关系”——这是考纲点名要理解的结论。这段证明把三个世界焊在一起：子式（行列式）、无关性（线性组合）、方程组的解（克拉默法则）；看懂“第 ③ 步靠延长性质、第 ⑤ 步靠子式为零”这两处卡点，就抓住了子式观点与向量组观点互译的全部机制。附带的那条结论更是解题利器：“找极大无关组”与“找最高阶非零子式所在的行列”是同一件事，遇到“已知 $r(A)=2$，求某参数使某两列构成极大无关组”的题可直接定位。
\end{fdeep}

% ---------- 印刷 88–89（§4 推论 1、2、3）----------
\begin{fdeep}{初等变换下的秩、阶梯形与极大无关组（北大 三.4 推论 1--3，印刷 88--89）}
\begin{itemize}[leftmargin=2.2em,itemsep=0.25em]
\item \textbf{推论 1}：矩阵的初等列变换和初等行变换都不改变该矩阵的秩、列秩和行秩。
      只对列变换说明，行变换同理：$A$ 经初等列变换变成 $B$ 时，两者的列向量组等价，
      故列秩不变；再串成一条链“$A$ 的行秩 $=A$ 的秩 $=A$ 的列秩 $=B$ 的列秩 $=B$ 的秩 $=B$ 的行秩”，
      可见三者被同时保持（链中第二、第三个等号用的是定理 4）。
\item \textbf{推论 2}：$A$ 的秩等于 $A$ 在初等行变换下所得阶梯形矩阵中\textbf{非零行的数目}。
      理由：阶梯形的 $r$ 个非零行各取首个非零元素所在的列 $i_1,\dots,i_r$，交出的 $r$ 阶子式
      等于 $b_{1i_1}b_{2i_2}\cdots b_{ri_r}\ne0$；而任何 $r+1$ 阶子式至少含一行全为零，
      故最高阶非零子式恰为 $r$ 阶，秩就是 $r$。这就是“算秩一律化阶梯形、数非零行”的依据。
\item \textbf{推论 3}：设 $A$ 的阶梯形是上述矩阵 $B$，则 $B$ 的第 $i_1,\dots,i_r$（主元）列组成
      $B$ 的列向量组的一个极大线性无关组；因而 $A$ 的第 $i_1,\dots,i_r$ 列组成 $A$ 的列向量组的
      一个极大线性无关组。
\end{itemize}
【反哺点】服务考纲〔矩阵〕“掌握用初等变换求矩阵的秩和逆矩阵的方法”与〔向量〕“会求向量组的极大线性无关组”。三条推论就是考场上的标准动作：化阶梯形 $\to$ 数非零行得秩 $\to$ 取主元列得极大无关组。要点有二：① 初等变换保持秩，是因为它把行（列）向量组换成等价组，而等价组秩相同；② 只有推论 2 用的是“阶梯形”这一具体形状，其余都要靠“等价”这一步过渡，答卷时不要跳过“向量组等价故秩相同”这句话。
\end{fdeep}

% ---------- 印刷 89（§4 推论 4、定理 5）----------
\begin{fdeep}{方阵情形：$|A|=0$、行（列）向量组线性相关与齐次方程组有非零解（北大 三.4 推论 4、定理 5，印刷 89）}
\textbf{推论 4}：设 $A$ 为 $n$ 阶方阵，则 $A$ 的列向量组（行向量组）线性相关的充要条件是 $|A|=0$；
线性无关的充要条件是 $|A|\ne0$。证明：$A$ 的最高阶子式就是 $n$ 阶行列式 $|A|$，
而“列向量组线性相关 $\iff$ 秩 $<n\iff$ 最高阶子式 $|A|=0$”。
\textbf{定理 5}：齐次线性方程组有非零解的充要条件是它的系数矩阵的行列式 $|A|=0$；
只有零解的充要条件是 $|A|\ne0$。（逆用即得克拉默法则的逆定理：$n$ 个方程 $n$ 个未知量的
非齐次方程组有唯一解 $\iff|A|\ne0$。）
【反哺点】服务考纲〔矩阵〕“理解矩阵的秩的概念”与〔线性方程组〕“掌握非齐次线性方程组有解和无解的判定方法”。方阵情形下四句话等价：$|A|=0$、行（列）向量组线性相关、$r(A)<n$、齐次方程组有非零解。做“含参数方程组解的个数”题时第一刀就砍在 $|A|$ 上：$|A|\ne0$ 直接唯一解；$|A|=0$ 说明列向量组相关、秩掉下来了，接着才去定秩究竟是 $n-1$ 还是更小。这条链与第 3 章的“可逆等价刻画”重合的部分请主控去重。
\end{fdeep}

% ===========================================================
%  跳过页/段记录（按「只取反哺解题的深化知识」原则）
% -----------------------------------------------------------
%  beida p90（印刷 78）定义 9 线性组合、定义 10 线性表出与等价
%        —— 属基础概念，已在 9 讲正文覆盖，不重复提取。
%  beida p91（印刷 79）等价的三条性质（自反/对称/传递）、定义 11 线性相关
%        —— 同上，基础概念；结论已并入第 1 块（方程组判定）。
%  beida p92（印刷 80）定义 11 与 11' 一致性证明、定义 12 线性无关
%        —— 基础概念，跳过（"无关 ⟺ 只有零解"已并入第 1 块）。
%  beida p93-p94（印刷 81-82）例题：判断 α=(2,-1,3,1) 组是否相关 —— 例题跳过；
%        其中的一般性结论（方程组判定、延长性质、定理 2）已提取为第 1、2 块。
%  beida p95（印刷 83）定义 13 后的极大无关组算例 —— 例题跳过。
%  beida p96-p97（印刷 84-85）定义 15 的行/列向量组算例、例 1 计算 k 阶子式
%        —— 例题跳过（定义 15 本身是基础概念，正文已覆盖）。
%  beida p101（印刷 89）例 2 求极大线性无关组 —— 例题跳过；
%        定理 6（克拉默法则及其逆定理）—— 属克拉默法则而非秩理论，前章已收，
%        仅在第 8 块末尾以一句话提及。
%  beida p102 下半-p103（印刷 90 下半-91）"应用举例：平板在热平衡下的温度分布"
%        （网格局域网 9 点方程组、图 1、用最大值原理证解唯一）—— beida p103
%        全部为该应用实例，按深化提取原则跳过；p102 的定理 6 部分同上。
% ===========================================================
